% Following magic comments allow for compilation of root file
% !TEX root = ../../../../temp_manuscript.tex

% Each chapter can start with a designed divider page. Drop a PDF named
% IntroductionTitlePage.pdf (digital) / IntroductionTitlePageBleed.pdf (print)
% in this folder and uncomment the block below.
%\ifx\setbleed\empty
%	\includepdf[pages=1, pagecommand={\thispagestyle{empty}}, fitpaper=true]{IntroductionTitlePage.pdf}
%\else
%	\includepdf[pages=1, pagecommand={\thispagestyle{empty}}, fitpaper=true]{IntroductionTitlePageBleed.pdf}
%\fi
\chapter{Introduction}\label{chap: introduction}
\newpage
% Count floats per chapter
\counterwithin{figure}{chapter}
\counterwithin{table}{chapter}

% Opening words for the introduction.
Primary bone and soft-tissue tumors (BSTT) are among the rarest neoplasms in humans. 
They can arise at any age and at almost any anatomical site, originating from connective tissue cells in 
bone, muscles, fat, blood vessels and cartilage. 
These lesions span a wide spectrum, from benign tumors to malignant sarcomas, with variable clinical behavior and treatment approaches.
Their rarity, combined with diverse subtypes and characteristics, poses substantial challenges for diagnosis and prognosis. 
Effective management for BSTT therefore requires a multidisciplinary approach involving radiologists, pathologists, surgeons, and oncologists. 
Within this framework, pathological examination of tissue samples remains the gold standard for definitive diagnosis, 
while radiological imaging, by macroscopically visualizing the tumor and its surrounding tissue, 
plays a central role across all stages of BSTT management, 
from initial assessment and biopsy guidance to surgical planning and follow-up.

The imaging volumn continues to grow over years under the development of new acquisition techniques and new modalities. 
Consequently, the workload on for clinicians, such as radiologists, who work on medical imaging has increased considerablely over decades.
Artificial intelligence (AI) offers new opportunities to support these clinicians in many clinical tasks.

Under this background, this thesis aims to develop and evaluate AI-based tools to support the diagnosis and management of BSTT, 
with a particular focus on routinely acquired imaging modality such as magnetic resonance imaging (MRI). 


The rarety and hetrogeineity of BSTT generally make it difficult for clinicians to gain sufficient experience in their diagnosis and management.



\section{Clinical management of bone and soft-tissue tumor}
TODO: write the introduction, citing sources like this \autocite{exampleKey2029}.



\section{AI in healthcare}
TODO: outline referring to e.g. \autoref{chap: chapter01} and \autoref{chap: chapter02}.




\section{Research aims and thesis outline}
